\documentclass[article]{jss}

\usepackage{booktabs}
\usepackage{array}
\usepackage{tabularx}
\usepackage{amsmath}
\usepackage{amssymb}
\usepackage{lmodern}
\usepackage{float}
\usepackage{placeins}
\usepackage{xr}
\externaldocument[supp-]{faissR_jss_supplement}

\hypersetup{hypertexnames=false}
\setlength{\emergencystretch}{2em}

\newcommand{\class}[1]{\code{#1}}
\newcommand{\fct}[1]{\code{#1()}}
\newcommand{\method}[1]{\code{#1}}

\author{Moussa Kassim$^{1,2,\dagger}$
  \And Martin Ocharo$^{1,2,\dagger}$
  \And Dalia Ahmed$^{1}$
  \AND Dupe Ojo$^{1}$
  \And Alessia Vignoli$^{3,4}$
  \And Leonardo Tenori$^{3,4}$
  \AND Dinesh Gupta$^{5}$
  \And Silvano Piazza$^{6,7}$
  \And Stefano Cacciatore$^{1,2,\ast}$}
\Plainauthor{Moussa Kassim, Martin Ocharo, Dalia Ahmed, Dupe Ojo, Alessia Vignoli, Leonardo Tenori, Dinesh Gupta, Silvano Piazza, Stefano Cacciatore}

\title{\pkg{faissR}: Nearest-Neighbor Search with FAISS and cuVS in \proglang{R}}
\Plaintitle{faissR: Nearest-Neighbor Search with FAISS and cuVS in R}
\Shorttitle{FAISS and cuVS Nearest-Neighbor Search in \proglang{R}}

\Abstract{
Nearest-neighbor search is widely used in statistical learning, clustering,
and local prediction. We present \pkg{faissR}, a Python-free \proglang{R}
interface to the Facebook Artificial Intelligence Similarity Search (FAISS)
library and optional NVIDIA Compute Unified Device Architecture (CUDA)
support through RAPIDS cuVS. The package integrates these libraries rather
than proposing a new approximate nearest-neighbor algorithm. A common
interface supports exact and approximate search with Euclidean, cosine, and
correlation distances. The interface accepts dense \proglang{R} double
matrices and float32 matrices. FAISS and cuVS operate on float32 coordinates,
so double inputs are converted at the compiled boundary. Users can reuse fitted
indexes on the central processing unit (CPU) and retain results on a graphics
processing unit (GPU) for use by compiled \proglang{R} packages. Returned
metadata identify the search implementation, input conversion, index reuse,
and result location. We illustrate the interface and evaluate correctness,
recall, and runtime across nine datasets, including comparisons with seven
established \proglang{R} interfaces. Relative performance depended on the
dataset and query workload, with index construction and repeated queries
favoring different choices. An optional experimental selector uses calibration
on an NVIDIA L40S for searches that include index construction and query every
observation. Reproducible examples and benchmark code accompany the package.
}

\Keywords{nearest neighbors, approximate search, FAISS, CUDA, float32, \proglang{R}}
\Plainkeywords{nearest neighbors, approximate search, FAISS, CUDA, float32, R}

\Address{
  \textsuperscript{1}Bioinformatics Unit\\
  International Centre for Genetic Engineering and Biotechnology (ICGEB)\\
  Cape Town 7925, South Africa\\
  \textsuperscript{2}Department of Integrative Biomedical Sciences, Institute
  of Infectious Disease \& Molecular Medicine (IDM), University of Cape Town\\
  Cape Town 7925, South Africa\\
  \textsuperscript{3}Department of Chemistry ``Ugo Schiff'', University of
  Florence, Sesto Fiorentino, Italy\\
  \textsuperscript{4}Magnetic Resonance Center (CERM), University of Florence,
  Sesto Fiorentino, Italy\\
  \textsuperscript{5}Translational Bioinformatics Group, International Centre
  for Genetic Engineering and Biotechnology (ICGEB), New Delhi, India\\
  \textsuperscript{6}Computational Biology Group, International Centre for
  Genetic Engineering and Biotechnology (ICGEB), Trieste, Italy\\
  \textsuperscript{7}Bioinformatics Facility, Department of Cellular,
  Computational and Integrative Biology (CIBIO), University of Trento,
  Trento, Italy\\[1ex]
  \textsuperscript{\(\dagger\)}Moussa Kassim and Martin Ocharo contributed
  equally and share first authorship.\\
  \textsuperscript{*}Corresponding author:\\
  Stefano Cacciatore, \email{stefano.cacciatore@icgeb.org}\\
  Open Researcher and Contributor ID (ORCID):
  \url{https://orcid.org/0000-0001-7052-7156}\\
  Project: \url{https://github.com/tkcaccia/faissR}
}

\begin{document}

\section[Introduction]{Introduction} \label{sec:introduction}

For a reference matrix $X \in \mathbb{R}^{n \times p}$ and query matrix
$Q \in \mathbb{R}^{m \times p}$, nearest-neighbor search returns the $k$
reference rows closest to each query. The operation underlies local regression
and classification, clustering, anomaly detection, manifold learning, and
exploration of high-throughput biological measurements. Exhaustive self-search
has quadratic distance-evaluation cost, whereas approximate nearest-neighbor
(ANN) methods exchange some accuracy for lower time or memory.

Several compiled nearest-neighbor interfaces are available in the
\proglang{R} ecosystem \citep{R}. \fct{RANN::nn2} exposes exact or
error-bounded search
with ANN k-d and box-decomposition (BD) trees,
while \fct{FNN::get.knn} and \fct{FNN::get.knnx} provide k-d-tree, cover-tree,
cosine-distance search (\code{algorithm = "CR"}), and
linear brute-force searches \citep{RANN,FNN}. \fct{Rnanoflann::nn} provides a
nanoflann-based alternative \citep{Rnanoflann}. Approximate-search interfaces
include the Annoy index
classes in \pkg{RcppAnnoy}; \fct{hnsw\_knn}, \fct{hnsw\_build}, and
\fct{hnsw\_search} in \pkg{RcppHNSW}; and \fct{nnd\_knn},
\fct{rnnd\_build}, and \fct{rnnd\_query} in \pkg{rnndescent}
for nearest-neighbor descent (NN-descent)
\citep{RcppAnnoy,RcppHNSW,rnndescent,DongMosesLi2011}. Within Bioconductor,
\pkg{BiocNeighbors} provides \fct{findKNN}, \fct{queryKNN}, and
\fct{buildIndex} as a common interface to exhaustive, k-means
k-nearest-neighbor (kNN),
vantage-point-tree, Annoy, and hierarchical navigable small-world (HNSW)
implementations \citep{BiocNeighbors,MalkovYashunin2020}.
These packages provide several choices for CPU search. \pkg{faissR} adds
access to the broader FAISS index family and optional cuVS execution through
one native \proglang{R} interface. FAISS includes inverted-file (IVF),
product-quantization (PQ), and graph indexes, as well as exhaustive search
\citep{JohnsonDouzeJegou2021,DouzeEtAl2024}. The surveyed R interfaces return
ordinary objects in host memory, whereas \pkg{faissR} can retain neighbor
results in NVIDIA device memory for direct use by another compiled R package.
We compare the CPU interfaces and assess GPU-resident output separately
through \fct{nn\_gpu}.

Table~\ref{tab:r-ecosystem} compares the search families, metric scope, index
reuse, and CUDA support of the principal R interfaces. Several packages
support reusable CPU indexes; \pkg{faissR} combines reusable indexes with a
unified FAISS/cuVS CPU and CUDA interface.

\begin{table}[t]
\centering
\footnotesize
\setlength{\tabcolsep}{2.5pt}

\begin{tabularx}{\textwidth}{@{}
    >{\raggedright\arraybackslash}p{0.14\textwidth}
    >{\raggedright\arraybackslash}p{0.22\textwidth}
    >{\raggedright\arraybackslash}X
    >{\centering\arraybackslash}p{0.12\textwidth}
    >{\centering\arraybackslash}p{0.07\textwidth}@{}}
\toprule
Package & Principal search family & Public metric/scope summary & Reusable index & CUDA \\
\midrule
\pkg{RANN} & ANN k-d/BD trees & Euclidean; exact or error-bounded CPU call & No & No \\
\pkg{FNN} & k-d, cover, CR, linear & Exact CPU search through several public algorithms & No & No \\
\pkg{Rnanoflann} & nanoflann k-d tree & Exact or ANN CPU tree search & No & No \\
\pkg{RcppAnnoy} & Annoy forests & Approximate CPU index objects & Yes & No \\
\pkg{RcppHNSW} & hnswlib & Approximate HNSW CPU search & Yes & No \\
\pkg{rnndescent} & NN-descent & Approximate CPU graph construction/query & Yes & No \\
\pkg{BiocNeighbors} & Exact, tree, Annoy, HNSW & Common Bioconductor CPU interface & Yes & No \\
\pkg{faissR} & FAISS Flat/IVF/PQ; HNSW and cuVS & Exact and ANN CPU/CUDA under one metric contract & Yes & Yes \\
\bottomrule
\end{tabularx}
\caption{Feature-level comparison with principal R nearest-neighbor
interfaces. ``Reusable'' identifies a public application programming
interface (API) that can retain an index across queries within a session.}
\label{tab:r-ecosystem}
\end{table}

RAPIDS cuVS provides CUDA vector search and is integrated with recent FAISS
builds \citep{cuVS,FaissCuVS}. Connecting these libraries to R requires more
than exposing their search functions. R commonly stores double-precision
matrices by column, whereas the libraries predominantly use single-precision
(float32) arrays stored by row. The interface must also return consistently
ordered distances, support repeated queries without rebuilding compatible
indexes, and make transfers between host and device memory explicit.

\pkg{faissR} handles these operations in compiled code and accepts direct
\pkg{float} input where supported \citep{float}. It returns information about the
implementation and conversions used for each call. A C-callable interface
allows other packages to consume GPU results without first copying them to
the host. These integration features are the main contribution; automatic
method selection is an optional experimental facility. The following
sections explain the software and evaluate its correctness, recall, and
runtime against existing R implementations.

\FloatBarrier
\section{Statistical and software design} \label{sec:design}

\subsection{Distances and recall}

The public metrics are Euclidean distance,
\begin{equation}
d_E(x,q)=\left\|x-q\right\|_2,
\end{equation}
cosine distance,
\begin{equation}
d_C(x,q)=1-\frac{x^\top q}{\|x\|_2\|q\|_2},
\end{equation}
and correlation distance,
\begin{equation}
d_R(x,q)=1-\frac{(x-\bar{x}\mathbf{1})^\top
(q-\bar{q}\mathbf{1})}{\|x-\bar{x}\mathbf{1}\|_2
\|q-\bar{q}\mathbf{1}\|_2}.
\end{equation}
Returned Euclidean distances use the ordinary $L_2$ scale. When a library
returns squared distances, the compiled interface takes their square roots.
Cosine search therefore normalizes each row, while correlation search first
subtracts the row mean and then normalizes. These operations have two important
edge cases. Cosine distance cannot be calculated for an all-zero row because
that row has no direction. Correlation distance cannot be calculated for a
constant row because subtracting its mean produces an all-zero row.

\pkg{faissR} handles these rows differently on CPU and CUDA. On CPU, two
all-zero rows have cosine distance 0, while an all-zero row and a nonzero row
have cosine distance 1. Similarly, two constant rows have correlation distance
0, while a constant row and a nonconstant row have correlation distance 1.
These package-defined values allow CPU search to continue. For CPU Flat search,
ties created by this rule are ordered by row number. CUDA search stops with an
error and does not silently switch to the CPU. Before running a search,
\fct{nn\_metric\_preflight} can identify the affected rows in the reference and
query matrices and report whether the requested device can proceed. With
\code{backend = "auto"}, the answer depends on which device is selected. Search
results record the distance definition and ordering in their metadata; the
field-level contract is given in the supplement and reference manual.

All devices reject \code{NA}, \code{NaN}, and positive or negative infinity.
FAISS/cuVS routes operate on float32 coordinates. Ordinary R doubles are
converted at the compiled boundary, so neighbors close to a rounding or tie
boundary can differ from a double-precision exhaustive calculation. For
normalized float32 routes, means and squared norms are accumulated in double
precision and normalized coordinates are stored as float32; native double
routes center and normalize in double precision. Native exact and candidate
scorers order equal distances by one-based row identifier. External providers
may resolve ties differently. For example, when the $k$th and $(k+1)$th
observations are the same distance from a query, either observation is a valid
$k$th neighbor. Two exhaustive implementations can therefore return different
row identifiers even though their ordered distances agree. The exact-reference
audit treats corresponding distances $a$ (observed) and $b$ (reference) as
equal when
$|a-b|\leq 10^{-5}+10^{-4}|b|$. The first term allows for rounding near zero,
and the second allows a slightly larger difference for larger distances. This
tolerance is used only by the exact-reference audit. It does not change the
reported distances or determine whether an approximate method attains
\code{target\_recall}.

The \code{metric} argument accepts only \code{"euclidean"}, \code{"cosine"},
and \code{"correlation"}. Other distances and aliases are rejected.
The package checks that the requested method, metric, and device can be used
together; an unsupported explicit request produces an error without changing
any of them.
\code{nn\_capabilities(runtime = TRUE)} exposes this matrix together with local
FAISS, CUDA, and cuVS availability, allowing applications and benchmark scripts
to preflight a request.

For exact neighbor identifiers $N_i^{\star}(k)$ and candidate identifiers
$N_i(k)$, query-level recall is
\begin{equation}
r_i(k)=\frac{|N_i(k)\cap N_i^{\star}(k)|}{k},
\qquad
\bar r(k)=\frac{1}{m}\sum_{i=1}^{m}r_i(k).
\label{eq:recall}
\end{equation}
Throughout this article, \emph{target recall} is a query-set summary. In a
call with \code{target\_recall} $= \tau$, $\tau$ is the requested lower
threshold for mean recall@$k$, $\bar r(k)$, across the evaluated query rows.
It is not a requirement that every individual query attain $r_i(k)\geq\tau$.
The public argument accepts the three tiers 0.90, 0.95, and 0.99; it neither
rounds other values nor interpolates parameters. Because exact neighbors are
generally unavailable during an ordinary user call, the package does not
measure recall on the user's data. Instead, automatic tuning and method
selection use the calibration entry associated with the requested tier.

In calibration, an approximate configuration was eligible when its observed
mean recall reached $\tau$ on the calibration query sample. In validation,
target attainment required the same threshold for every independently sampled
query set. Repeated timings of one query set were execution checks rather than
additional independent recall samples. Identifier overlap can penalize an
equally valid neighbor with a different identifier at a tied boundary. The
reported approximate-search results use this unadjusted measure and no
bootstrap confidence bound. Exact methods are assessed separately by comparing
identifiers or accepting equivalent sorted distances under the tolerance
above; we label methods passing these checks \emph{exact-audited}. For paired
comparisons, both methods must pass the specified threshold, although their
recall values may differ. Section~\ref{sec:evaluation} describes query sampling
and the scope of inference.

\subsection{Methods and providers}

Table~\ref{tab:methods} summarizes the search methods. A method name describes
the requested strategy, while returned metadata identify the implementation
that executed it. The experimental CPU methods \method{nsg\_style},
\method{vamana\_style}, and \method{nndescent\_style} are not complete
implementations of the published navigating spreading-out graph (NSG), Vamana,
or NN-descent algorithms
\citep{FuEtAl2019,SubramanyaEtAl2019,DongMosesLi2011}. Published NSG and Vamana
construct navigable graph indexes and define how new queries traverse them;
NN-descent iteratively constructs an approximate $k$-nearest-neighbor graph.
The package-owned routes instead refine a bounded candidate list for each row
of a self-search and rescore the retained candidates exactly. They retain,
respectively, an NSG-like relative-neighborhood rule, Vamana's $\alpha$-scaled
robust-pruning rule, or bounded neighbor and reverse-neighbor updates. They
omit the complete index-construction and query-traversal procedures, hence the
\code{\_style} suffix. On CUDA, \method{nndescent\_style} can instead use the
external cuVS implementation; returned metadata identify which route ran.
Supplementary Table~\ref{supp-tab:supp-derived-methods} and the accompanying
pseudocode give a stage-by-stage comparison.

The names \method{exact}, \method{flat}, and \method{bruteforce} all request
exhaustive search but differ in how they choose a library. Use \method{exact}
for the preferred available implementation, \method{flat} to require a FAISS
Flat index, and \method{bruteforce} to prefer direct cuVS brute force on CUDA.
On CPU, \method{bruteforce} uses FAISS Flat. Several names may therefore use
the same implementation on a given installation. Both the requested name and
the actual \code{backend\_used} are retained. The supplement gives the full
order of preference and capability requirements.

CAGRA is the official cuVS name for its highly parallel graph-construction
and approximate-search method.

\begin{table}[t]
\centering

\small
\begin{tabularx}{\textwidth}{@{}>{\raggedright\arraybackslash}p{0.27\textwidth}X@{}}
\toprule
Public method & Family, status, and principal controls \\
\midrule
\method{exact}, \method{flat}, \method{bruteforce} & Exhaustive search; query batch and resource reuse. \\
\method{grid} & Exact low-dimensional search; spatial cell width and query expansion. \\
\method{hnsw} & Navigable graph; graph degree, construction, and search effort. \\
\method{ivf} & Inverted file; number of lists and probed lists. \\
\method{ivfpq}, \method{ivfpq\_fastscan} & Provider-backed quantization; lists, probes, subquantizers, code width, block size, and refinement. \\
\method{cagra} & CAGRA graph search; graph degrees and search width. \\
\method{nndescent\_style} & Experimental package route with projection seeding and bounded, fixed-iteration graph updates; an external cuVS implementation may be used instead. \\
\method{nsg\_style}, \method{vamana\_style} & Experimental package routes that apply, respectively, an NSG-like relative-neighborhood rule or Vamana's $\alpha$-scaled robust-pruning rule to an existing candidate list. \\
\bottomrule
\end{tabularx}
\caption{Nearest-neighbor method families exposed by \pkg{faissR} and their
principal controls.}
\label{tab:methods}
\end{table}

The supplement describes the experimental graph refinements, including
pseudocode, initialization, time and memory bounds, and differences from the
canonical algorithms (Supplementary Table~\ref{supp-tab:supp-derived-methods}).
Their correctness tests and calibration results are reported, but these
methods are excluded from the principal performance comparisons. The
low-dimensional grid method is checked separately as an exact method. IVF-PQ
and FastScan call the corresponding library implementations; their evidence
here covers route correctness and calibration.

With \code{method = "auto"}, the package chooses a method family. Once a method
is known, \code{tuning = "auto"} retrieves its calibrated parameter settings.
Both choices use deterministic lookups. Package-owned graph construction uses
the recorded
\code{faissR.approx\_knn\_seed} option and deterministic tie
ordering. FAISS/cuVS construction can contain provider-controlled randomized
or parallel stages; the same input, route parameters, seed, thread controls,
provider versions, and hardware are therefore required for procedural
reproduction, but GPU execution is not claimed to be bitwise deterministic.
The supplement lists the corresponding metadata fields.

\subsection{Data movement and fitted indexes}

Dense base-\proglang{R} matrices are converted once to contiguous float32
storage at the compiled boundary, whereas \pkg{float} matrices can avoid this
conversion. The \fct{knn} and \fct{predict} interfaces retain compatible fitted
indexes, separating index construction from repeated queries. On CUDA,
\fct{nn\_gpu} can return neighbor buffers that remain device-resident;
\fct{gpu\_knn\_to\_host} performs an explicit transfer to ordinary
\proglang{R} matrices. Returned metadata record the input representation,
selected provider, index reuse, result residency, and host-transfer count.
Detailed ownership, synchronization, and object-lifetime contracts are
provided in the supplement.

\section[The R interface]{The \proglang{R} interface} \label{sec:interface}

Table~\ref{tab:api} lists the main functions and their return objects.
We focus on nearest-neighbor search. Prediction by kNN and k-means clustering
are secondary interfaces, documented with examples in the package vignette.

The argument \code{backend} selects the execution device: CPU, CUDA, or
automatic device selection. The metadata field \code{backend\_used} identifies
the implementation, such as \code{faiss\_hnsw}. We use \emph{device} and
\emph{implementation} to distinguish these meanings in the text.

Automatic method selection is experimental. The CUDA policy was calibrated
on an NVIDIA L40S for full self-search, including index construction.
Section~\ref{sec:tuning} explains method selection and parameter tuning.

\begin{table}[t]
\centering

\small
\begin{tabularx}{\textwidth}{@{}>{\raggedright\arraybackslash}p{0.27\textwidth}>{\raggedright\arraybackslash}p{0.25\textwidth}X@{}}
\toprule
Interface & Main use & Result or lifecycle \\
\midrule
\fct{nn} & Host-returned neighbor search & One-based indices and distances in two $m\times k$ matrices. \\
\fct{nn\_gpu}, \fct{gpu\_knn\_to\_host} & Device-resident search and explicit transfer & An owning GPU result; conversion to host matrices occurs only when requested. \\
\fct{knn}, \fct{predict} & Classification or regression with repeated queries & Compatible fitted indexes are retained within the R session. \\
\fct{backend\_info}, \fct{nn\_capabilities} & Capability inspection & Available providers and supported method--device--metric combinations. \\
\fct{nn\_metric\_preflight} & Input diagnostics & Locations of non-finite or undefined rows and the requested device action. \\
\bottomrule
\end{tabularx}
\caption{Principal public interfaces. Detailed object and ownership contracts
are given in the supplement and reference manual.}
\label{tab:api}
\end{table}

\subsection{Choosing a method and preparing the data}

Exact search is a reference choice and is often competitive for a one-off call.
HNSW is a useful CPU starting point for repeated queries; IVF offers a
training/search tradeoff on large collections, and CAGRA provides CUDA graph
search. Table~\ref{tab:methods} summarizes the main controls and route status.
Experimental
automatic selection should be used within its documented calibration scope.

Every host result has one row per query and one column per returned neighbor.
Duplicate observations remain distinct row identifiers, so a non-self
duplicate may correctly have distance zero. For self-search,
\code{exclude\_self = TRUE} removes only the query row identifier; if an
approximate candidate list omitted that identifier, the first $k$ non-self
candidates are retained rather than discarding another row.

Inputs must be dense numeric matrices, data frames, or \pkg{float} matrices,
with observations in rows. Sparse, delayed, file-backed, and other matrix-like
objects are rejected before conversion. For single-cell workflows, retain
the sparse assay during preprocessing and supply a dense representation of
manageable size, such as principal component analysis (PCA) coordinates.
FAISS/cuVS searches convert R doubles to contiguous row-major float32 storage,
requiring approximately $4np$ additional bytes while the source matrix remains
in memory. Indexes and working space require further memory. Direct
\pkg{float} adapters can avoid numeric conversion, although a layout copy
may still be needed.

Native indexes are retained within the R session and released when their
owning objects are garbage-collected. Saving a model with \fct{saveRDS}
preserves its R components and training data. The serialized object excludes
the native index, so \fct{predict} rebuilds it after reload and reports
\code{query\_source = "nn"} for that call. GPU results cannot be serialized
or shared between processes. Both host and CUDA identifiers are signed
32-bit integers, limiting the interface to $2^{31}-1$ reference rows;
available memory and library limits usually impose smaller limits.
The supplement gives the detailed object-lifetime rules.
The package-wide backend can be set with \fct{faissR\_backend}, an
\proglang{R} option, or an environment variable. An explicit function argument
has highest precedence. An unsupported backend produces an error without
changing the requested device.

\subsection{Installation and hardware portability}

Functional CPU builds require \proglang{C++20}, a Fortran compiler, and a
toolchain-compatible FAISS installation. \pkg{faissR} discovers FAISS through
\code{pkg-config}, standard prefixes, or \code{FAISS\_HOME}; CUDA/cuVS support
is optional for CPU-only installation and becomes a fatal configure requirement
only when explicitly requested. Linux and macOS source builds are supported.
Native Windows CPU use requires an Rtools-compatible FAISS build; Windows
Subsystem for Linux 2 (WSL2) is the
documented practical route for NVIDIA libraries. Automated binary builders
without compatible system libraries can produce diagnostic-only builds that
load and report the missing library but cannot execute FAISS search.
The package does not vendor FAISS/cuVS and requires compiler and external-library
application binary interface (ABI) compatibility. Its downstream C-callable
interface is independently versioned as ABI 1 and checked by consumers before
use.

A typical macOS CPU installation starts with the following shell command:
\begin{CodeChunk}
\begin{CodeInput}
brew install faiss libomp
\end{CodeInput}
\end{CodeChunk}
Then install and load the package in R:
\begin{CodeChunk}
\begin{CodeInput}
R> install.packages("faissR_0.99.44.tar.gz", repos = NULL, type = "source")
R> library("faissR")
R> backend_info()
R> subset(nn_capabilities(runtime = TRUE), runtime_available)
\end{CodeInput}
\end{CodeChunk}
For Ubuntu releases without \code{libfaiss-dev}, the installation vignette gives
an exact FAISS 1.14.3 source build with GPU, Python, and tests disabled, followed
by a strict package installation using \code{FAISS\_HOME=/opt/faiss} and
\code{FAISSR\_REQUIRE\_FAISS=1}. A strict CUDA/cuVS build additionally sets
\code{CUDA\_HOME}, \code{CUVS\_HOME}, \code{FAISSR\_REQUIRE\_CUDA=1}, and
\code{FAISSR\_REQUIRE\_CUVS=1}; a missing requested provider is then fatal.
After installation, \fct{faiss\_available}, \fct{backend\_info}, and
\fct{nn\_capabilities} are the authoritative checks, followed by a small exact
CPU search. Users should not infer CUDA availability merely from the presence
of an NVIDIA device.

Tested environments and their compiler, FAISS, CUDA, and cuVS versions are
listed in Supplementary Table~\ref{supp-tab:install-matrix}; execution identities and check records are
retained in the replication archive.

\pkg{faissR} is MIT-licensed. It neither vendors nor redistributes FAISS
(MIT), cuVS (Apache-2.0), or the CUDA toolkit/driver (NVIDIA CUDA software
development kit (SDK) terms).
Those components are supplied by the user's system; the installed third-party
license record links the governing texts and states this distribution boundary.

\subsection{Diagnosing inputs and planning memory}

Capability and metric checks can be performed without launching a search:
\begin{CodeChunk}
\begin{CodeInput}
R> caps <- nn_capabilities(runtime = TRUE)
R> subset(caps, method %in% c("hnsw", "ivf", "cagra") &
+              metric == "euclidean",
+        c(method, backend, runtime_available, runtime_reason))
R> x_bad <- rbind(c(0, 0), c(1, 0), c(NA, 1))
R> nn_metric_preflight(x_bad, metric = "cosine", backend = "cuda")
\end{CodeInput}
\end{CodeChunk}
The preflight reports one-based non-finite and zero/constant-row locations and
whether the requested backend would accept them. The requested backend remains
unchanged. Before a large search, users should budget for the source R
matrix, approximately $4np$ bytes of contiguous float32 provider input,
provider index/workspace, and at least $mk(4+8)$ bytes for host integer
identifiers and double distances. \pkg{float} input can avoid retaining both
double and converted input representations. The current \method{auto} policy
has no memory-limit input, so a constrained workflow should choose an explicit route,
use a representative pilot subset, query in application-level batches, and
monitor host or device memory externally. The package reports failures rather
than silently moving a CUDA request to CPU.

On a different or unidentified GPU, \method{auto} remains available but emits one
warning per runtime model and records
\code{hardware\_extrapolated\_unvalidated}. Explicit methods retain the user's
choice without this warning. \code{tuning = "pilot"} and \code{"cache"}
tune parameters within the selected method; the method family remains fixed.

\subsection{A practical CPU workflow}

The Bioconductor \pkg{Biobase} package supplies \code{sample.ExpressionSet},
with measurements for 500 probes in 26 samples \citep{HuberEtAl2015}. We first
treat probes as observations and find probes with similar expression profiles.
We use this compact dataset to explain the interface; Section~\ref{sec:evaluation}
evaluates performance on substantially larger matrices.
\begin{CodeChunk}
\begin{CodeInput}
R> library("faissR")
R> data("sample.ExpressionSet", package = "Biobase")
R> expression_data <- Biobase::exprs(sample.ExpressionSet)
R> x <- scale(expression_data)
R> set.seed(1)
R> exact <- nn(x, k = 15, exclude_self = TRUE, backend = "cpu",
+             method = "exact", n_threads = 1)
\end{CodeInput}
\end{CodeChunk}
The columns are standardized by \fct{scale}; each row remains a probe. To
compare approximate results with these exact neighbors, we calculate the
mean fraction of shared identifiers and evaluate the three supported tiers:
\begin{CodeChunk}
\begin{CodeInput}
R> recall_at_k <- function(reference, candidate) {
+    mean(vapply(seq_len(nrow(reference)), function(i) {
+      length(intersect(reference[i, ], candidate[i, ])) / ncol(reference)
+    }, numeric(1)))
+  }
R> hnsw_results <- do.call(rbind, lapply(c(0.90, 0.95, 0.99), function(tier) {
+    candidate <- nn(x, k = 15, exclude_self = TRUE, backend = "cpu",
+                    method = "hnsw", tuning = "auto",
+                    target_recall = tier, n_threads = 1)
+    data.frame(requested_tier = tier,
+               observed_recall = recall_at_k(exact$indices, candidate$indices),
+               provider = attr(candidate, "backend_used"))
+  }))
\end{CodeInput}
\end{CodeChunk}

For index reuse, we turn to a second task: predicting sample type. Samples
now form the rows, so we transpose the expression matrix. We hold out every
fourth sample and select the 128 most variable probes using training data
only. The same training means and scales are then applied to the test samples:
\begin{CodeChunk}
\begin{CodeInput}
R> y <- droplevels(Biobase::pData(sample.ExpressionSet)$type)
R> test <- seq.int(4L, ncol(expression_data), by = 4L)
R> train <- setdiff(seq_len(ncol(expression_data)), test)
R> variance <- apply(expression_data[, train, drop = FALSE], 1L, var)
R> keep <- head(order(variance, decreasing = TRUE), 128L)
R> training <- scale(t(expression_data[keep, train, drop = FALSE]))
R> queries <- scale(t(expression_data[keep, test, drop = FALSE]),
+                  center = attr(training, "scaled:center"),
+                  scale = attr(training, "scaled:scale"))
\end{CodeInput}
\end{CodeChunk}
We can now time construction and repeated predictions, then check that the
fitted index was reused:
\begin{CodeChunk}
\begin{CodeInput}
R> build_sec <- system.time(model <- knn(
+    training, y[train], k = 5, backend = "cpu", method = "hnsw",
+    n_threads = 1))["elapsed"]
R> first_sec <- system.time(prediction <- predict(model, queries))["elapsed"]
R> times <- replicate(3, system.time(
+    prediction <- predict(model, queries))["elapsed"])
R> attr(prediction, "faissR_nn")$query_source
\end{CodeInput}
\begin{CodeOutput}
[1] "fitted_index"
\end{CodeOutput}
\end{CodeChunk}
The companion \code{practical\_cpu\_example.R} prints the recall results and
summarizes these timings. Representative output is:
\begin{CodeChunk}
\begin{CodeOutput}
 requested_tier observed_recall   provider
           0.90           0.972 faiss_hnsw
           0.95           0.991 faiss_hnsw
           0.99           0.998 faiss_hnsw
 build_sec first_sec median_reused_sec query_source
     0.002     0.002             0.001 fitted_index
\end{CodeOutput}
\end{CodeChunk}

Here, \code{faiss\_hnsw} identifies the implementation and
\code{fitted\_index} confirms that prediction reused a native index.
The observed recall is measured on this example; the nested
\code{approximation} metadata describe the calibration used to choose
parameters. The millisecond timings illustrate the calls and depend on the
machine and timer resolution. Controlled performance comparisons use the
experiments described in Section~\ref{sec:evaluation}.

\subsection{GPU-resident output}

When CUDA support is compiled, the nearest-neighbor result can remain on the
GPU:
\begin{CodeChunk}
\begin{CodeInput}
R> g <- nn_gpu(x, k = 30, exclude_self = TRUE,
+            method = "auto", metric = "euclidean",
+            tuning = "auto", target_recall = 0.99)
R> c(residency = g$result_residency,
+    host_copies = g$device_to_host_result_copies)
\end{CodeInput}
\begin{CodeOutput}
   residency host_copies
      "cuda"         "0"
\end{CodeOutput}
\begin{CodeInput}
R> host_result <- gpu_knn_to_host(g)
\end{CodeInput}
\end{CodeChunk}
The opaque object is intended for downstream compiled packages; direct
inspection in \proglang{R} requires the explicit transfer. The object reports
its device, layout, element types, owner, synchronization contract, and whether
serialization or interprocess sharing is supported. A consumer selects the
recorded device, may use its own CUDA stream after the synchronized handoff,
and must complete that work before releasing the owning handle.

A minimal compiled consumer is supplied in the replication repository under
\path{validation/gpu_resident_interoperability/consumer/faissRGpuConsumer}.
It checks the C-callable ABI, retains the owning handle, selects the producing
device, and synchronizes its work before releasing the owner. The supplement
documents buffer layout and lifetime. Quantitative transfer performance lies
outside this interface example.

\subsection{Automatic method selection and tuning} \label{sec:tuning}

Three arguments control different decisions. \code{backend = "auto"} resolves
the execution device, \code{method = "auto"} invokes the experimental
method-family policy, and \code{tuning = "auto"} selects parameters within the
resolved method. An unsupported explicit request fails before execution.
The requested \code{target\_recall} is one of 0.90, 0.95, or 0.99.

For each dataset, device, method, metric, and value of $k$, we evaluated a
predefined grid of parameter configurations. Each run included preprocessing,
index construction or training, and a search for every observation against
the dataset. We call this a \emph{cold self-search}: no fitted index is reused.
We then measured recall against exact neighbors for $\min(n,1024)$ sampled
calibration rows. A configuration was considered complete when both a
successful timing and a recall record were available. Unsupported settings,
failures, and runs stopped by time or memory limits were recorded but could
not be selected.

For an approximate method, we selected the fastest completed configuration
whose mean recall reached the requested tier $\tau$. Ties were resolved by
higher recall and then configuration identifier. Exact methods qualified
through separate exactness checks. If no approximate configuration reached
$\tau$, we retained the one with the highest recall and marked it as below
target, using runtime to break recall ties. Each configuration was timed
once, so selection used its observed runtime rather than an average over
repeated measurements.

These recommendations supply the lookup used by \code{tuning = "auto"}.
The separate \code{method = "auto"} policy chooses a method family before
applying its stored parameters. Selection rules use data shape, query count,
self-search or external-query mode, metric, $k$, tier, and available libraries.
The current policy omits expected index reuse and a memory budget. The complete
grid, missing results, timing sensitivity, and changes from calibration to
validation are reported in Supplementary Tables~\ref{supp-tab:grid},
\ref{supp-tab:calibration-missing}, \ref{supp-tab:calibration-stability}, and
\ref{supp-tab:calibration-validation}.

\section{Evaluation} \label{sec:evaluation}

\subsection{Datasets and query sampling}

We used nine datasets with $873$ to $5.22$ million observations and $11$ to
$16,384$ variables (Table~\ref{tab:datasets}). The image datasets comprise
Columbia Object Image Library (COIL-20), Fashion-MNIST, Modified National
Institute of Standards and Technology (MNIST), United States Postal Service
(USPS), and ImageNet features
\citep{Nene1996,Xiao2017,LeCun1998,YouShung2022,Hull1994,
Russakovsky2015,Oquab2024}. The cytometry datasets are FR-FCM-ZYRM, flow18,
and mass41 \citep{VanUnen2017,Belkina2019}; MetRef contains metabolomics data
\citep{Cacciatore2014}. Supplementary Table~\ref{supp-tab:dataset-provenance}
gives their sources, preprocessing, and redistribution conditions. Labels
were not used in search, tuning, or recall calculations. Calibration,
exact references, and CUDA selector validation used float32 input. No dataset
underwent additional dimension reduction within the benchmark.

We timed full self-searches ($Q=X$, $m=n$), excluding each observation from
its own neighbors. Timings included preprocessing, index construction, the
complete search, and output creation. We assessed recall on up to 1,024
distinct rows sampled without replacement, using exact references at $k=100$
and retaining the first $k$ neighbors for $k\in\{15,30,50,100\}$. Datasets
with at most 1,024 rows used all observations, so MetRef contributed 873
queries. The requested tiers were $\{0.90,0.95,0.99\}$.

Two validation seeds defined the independent recall samples. Each search was
repeated three times with the query sample held fixed within a seed. Recall was
recomputed after each run as an execution check, and all three runs for both
seeds had to meet the tier. Timings were summarized across repeats within seed.
Calibration and validation sampled rows separately from the same datasets, so
this design tests replication within a dataset.

We also examined leave-one-dataset-out (LOODO) selection. All measurements
from the omitted dataset were removed before choosing a method from the
remaining datasets. Selection used the same predefined shape group when
available, otherwise the three nearest training datasets in
$\log(n)$--$\log(p)$ space. The omitted dataset was used only for evaluation.
This post hoc analysis used a different set of methods from the installed
policy and is reported separately as a sensitivity analysis.

The exact references comprise $9\times3\times3=81$ real-data records
(dataset $\times$ metric $\times$ seed) and $2\times1\times3=6$ Euclidean
synthetic-spatial records, for 87 records in total
(Supplementary Table~\ref{supp-tab:reference-key}).
The file \path{reference_record_dimensions.csv} records their dimensions,
query counts, seeds, fingerprints, implementations, and independent CPU
checks.

\begin{table}[t]
\centering

\begin{tabular}{@{}lrr@{}}
\toprule
Dataset & Rows & Variables \\
\midrule
COIL-20 & 1,440 & 16,384 \\
Fashion-MNIST & 70,000 & 784 \\
flow18 & 1,000,021 & 11 \\
FR-FCM-ZYRM & 5,220,347 & 32 \\
ImageNet features & 1,281,167 & 1,024 \\
mass41 & 965,282 & 14 \\
MetRef & 873 & 375 \\
MNIST & 70,000 & 784 \\
USPS & 11,000 & 256 \\
\bottomrule
\end{tabular}
\caption{Datasets used in the evaluation.}
\label{tab:datasets}
\end{table}

\subsection{Computing environment}

CPU jobs requested 12 threads on the heterogeneous University of Cape Town
(UCT) \emph{ada} partition. Each CPU pair ran on the same node. CUDA jobs used
one NVIDIA L40S with 46,068 mebibytes (MiB), compute capability 8.9, and driver
595.58.03. Both used \proglang{R} 4.5.3 on Debian 13, OpenBLAS 0.3.33, and a
checksummed Singularity image \citep{Singularity}. Supplementary
Table~\ref{supp-tab:supp-environment}
lists the software and hardware details.

The article accompanies \pkg{faissR} 0.99.44. The replication archive retains
the package, library, container, and data identifiers recorded for each
experiment. Runtime ratios compare calls within the same experiment.

\subsection{Comparisons and timing}

Correctness tests checked dimensions, finite and sorted distances, the
implementation used, metric invariances, undefined rows, and GPU residency.
Exact methods also underwent the tie-aware checks in Section~\ref{sec:design}.
Calibration evaluated all available methods. CUDA validation covered
\method{auto}, explicit Flat, brute force, and CAGRA.

The CPU studies addressed two questions: how public interfaces performed with
prespecified settings, and how HNSW implementations compared when tuned
independently. The fixed-configuration HNSW study compared \pkg{faissR} with
\pkg{BiocNeighbors} and \pkg{RcppHNSW} for Euclidean and cosine distance.
The latter two used graph degree 16, construction effort 200, and search
effort $\max(50,3k)$ through their respective public arguments. Correlation
was evaluated for \pkg{faissR} but was not part of these external pairs.
The supplement gives the full parameter names and settings, with dataset-level
results in Supplementary Table~\ref{supp-tab:supp-cpu-paired} and
Figure~\ref{supp-fig:supp-paired-cpu}.

For each matched comparison $c$, the speed ratio is
\begin{equation}
S_c=\frac{T_{c,\mathrm{second\ method}}}{T_{c,\mathrm{first\ method}}},
\label{eq:speed-ratio}
\end{equation}
so $S_c>1$ favors the first method named. We included pairs only when both
calls completed and met the specified mean-recall threshold. Ratios were
summarized within datasets before calculating the across-dataset median and
interquartile range (IQR). This treats datasets as the primary units and
avoids pooling absolute times from different problems.

The controlled HNSW comparison ran each pair on the same node, randomized
the initial method, and alternated order across five repeats. Each call ran
in a fresh R worker after data loading and an untimed 128-row warm-up, with
the fitted-index cache cleared before the cold call. Both methods received
the same double matrix, with conversion and output creation included in the
timer. Twelve threads were requested through the interfaces and environment
variables. We recorded the requested limits but did not instrument each
provider's internal thread utilization. Recorded node, allocation, and
execution-order fields were checked before pairing. CUDA timers ended after
synchronization.

The independently tuned HNSW experiment selected a configuration for each
implementation using calibration seed 20260706 and evaluated it using seed
20260807. It measured cold calls, direct index construction, and fitted-index
queries separately over five validation repetitions.

The broader comparison included \pkg{FNN}, \pkg{RANN}, \pkg{Rnanoflann},
\pkg{rnndescent}, \pkg{BiocNeighbors}, \pkg{RcppAnnoy}, and \pkg{RcppHNSW}.
We used only metrics supported by each interface: Euclidean for \pkg{FNN},
\pkg{RANN}, and \pkg{Rnanoflann}; Euclidean and cosine for the HNSW and Annoy
interfaces; and Euclidean, cosine, and correlation for \pkg{rnndescent}.
We distinguished three interpretations. Exact interfaces were paired with
\code{exact}, and HNSW interfaces with \code{hnsw}; these are same-family
comparisons. Because \pkg{faissR} has no Annoy implementation, the Annoy
interfaces were paired with \code{method = "auto"} as task-level
alternatives rather than as algorithm comparisons. Finally,
\pkg{rnndescent} was paired with the experimental package-owned
\code{nndescent\_style} route. That result is descriptive, appears only in
the supplement, and is excluded from the principal performance table. Pairs
shared the dataset, metric, $k$, validation seed, repeat, Slurm allocation,
and node; order was randomized. The complete results appear in Supplementary
Figure~\ref{supp-fig:supp-comprehensive-r}.

A separate workload experiment used 1, 32, and up to 1,024 external query rows
on five datasets. It recorded a cold call and repeated identical queries.
Here, construction time was estimated as
$\max(0,T_{\mathrm{cold}}-T_{\mathrm{warm}})$, rather than measured directly.
The total for $A$ query batches was estimated as
$T_{\mathrm{cold}}+(A-1)T_{\mathrm{warm}}$. The supplement describes the
reuse checks and the comparison with repeated one-shot Flat calls.

All timed CUDA comparisons used \fct{nn}, so every method returned matrices
in host memory. The GPU-resident assessment covered ownership, lifetime, and
explicit transfer; it did not include a transfer-performance benchmark.

\section{Results} \label{sec:results}

\subsection{Correctness and calibration}

All 48 CPU and 52 CUDA correctness test cases passed. The exact-reference
archive contained all 87 records and passed the identity and independent
CPU checks.

Calibration produced 6,453 of 6,804 planned combinations of method, dataset,
$k$, and recall target. The 351 unavailable combinations comprised 318 CPU
and 33 CUDA cases: 306 reached the 48-hour limit and 45 exhausted memory.
They were concentrated in FR-FCM-ZYRM (201), ImageNet features (144), and
flow18 (6). Supplementary Table~\ref{supp-tab:calibration-missing} gives the
method breakdown, with metric totals in the accompanying text.
Of 4,725 completed approximate-search recommendations, 3,653 (77.3\%)
reached the requested tier and 1,072 remained below target. A further 1,728
exhaustive recommendations were exact-audited.

Each configuration was timed once. The available timing-sensitivity diagnostic
is therefore the gap between the fastest and second-fastest eligible
configurations. Supplementary
Tables~\ref{supp-tab:calibration-stability} and
\ref{supp-tab:calibration-validation} report these gaps and the changes from
calibration to validation. Recall transferred more reliably than single-run
timing, particularly for CUDA IVF.

\subsection{CPU HNSW performance and index reuse}

We first compared separately tuned HNSW implementations. All 720 executions
completed. Of the 36 dataset--$k$ combinations evaluated for each comparator,
both implementations reached mean recall of at least 0.99 in 31 combinations
for \pkg{BiocNeighbors} and 32 for \pkg{RcppHNSW}. Only these combinations
contributed to the time ratios in Table~\ref{tab:tuned-hnsw}.

\pkg{faissR} had lower median cold-call and index-construction times than both
comparators. Relative to \pkg{BiocNeighbors}, the median ratios were 6.55 for
the cold call and 11.50 for index construction. Relative to \pkg{RcppHNSW},
they were 1.67 and 1.63. Once an index had been constructed, query times were
similar: \pkg{BiocNeighbors} took 0.90 times the \pkg{faissR} time, whereas
\pkg{RcppHNSW} took 1.09 times the \pkg{faissR} time. The largest differences
therefore occurred during index construction rather than during queries on an
existing index.

\begin{table}[H]
\centering

\begin{tabular}{@{}lrrrr@{}}
\toprule
Comparator & Eligible & Cold call & Index build & Fitted query \\
\midrule
\pkg{BiocNeighbors} & 31/36 & 6.55 & 11.50 & 0.90 \\
\pkg{RcppHNSW} & 32/36 & 1.67 & 1.63 & 1.09 \\
\bottomrule
\end{tabular}
\caption{Independently tuned CPU HNSW comparison. Eligible gives the number
of dataset--$k$ combinations, out of 36, in which both implementations reached
mean recall of at least 0.99 in every validation repeat. Time entries are
medians of $T_{\mathrm{comparator}}/T_{\mathrm{faissR}}$ across eligible
combinations; values above one favor \pkg{faissR}.}
\label{tab:tuned-hnsw}
\end{table}

A separate workload experiment examined repeated queries. All 60 CPU and 60
CUDA cases completed. On the CPU, querying an existing HNSW index was faster
than rebuilding a Flat index in each of the 15 HNSW cases. After including the
initial HNSW construction cost, 13 cases reached break-even within the
evaluated range; their median break-even point was 13 query batches. Because
the Flat baseline was rebuilt for every batch, this result compares HNSW reuse
with repeated one-shot exhaustive calls; it is
not a comparison with a reusable exact index. The
calculation and assumptions are given in Supplementary
Section~\ref{supp-sec:supp-workload}. The CUDA runs recorded no fitted-index
cache hit and therefore do not provide evidence about GPU index reuse.

\subsection{Comparison with existing R interfaces}

The seven-package experiment comprised 216 tasks and 6,480 recorded route
executions.
Of 4,104 external-versus-\pkg{faissR} pairs, 1,636 completed and met the
requested mean recall for both methods; 1,385 calls reached the time limit and
22 failed for other reasons. Table~\ref{tab:comprehensive-r} separates the
principal same-family comparisons from the Annoy-versus-\code{auto}
task-level comparisons. The experimental \pkg{rnndescent} comparison is
reported only in the supplement. Supplementary
Figure~\ref{supp-fig:supp-comprehensive-r} shows the class-level ratios and
their interquartile intervals together with the individual dataset medians.
Supplementary Table~\ref{supp-tab:supp-comprehensive-datasets} gives the
numerical dataset-specific values and eligible-pair counts.
Supplementary Table~\ref{supp-tab:supp-evidence-audit} summarizes test and
experiment completion across the evidence streams used in the article.
Per-dataset results for the fixed-configuration HNSW experiment appear in
Supplementary Table~\ref{supp-tab:supp-cpu-paired} and
Figure~\ref{supp-fig:supp-paired-cpu}.

\begin{table}[H]
\centering
\scriptsize

\begin{tabular}{@{}llrrrr@{}}
\toprule
Comparator & \pkg{faissR} method & Datasets & Matched/planned & Timeouts & Median [IQR] \\
\midrule
\multicolumn{6}{@{}l}{\textit{Panel A: same-family comparisons}} \\
\addlinespace[0.25em]
\pkg{FNN} exact & \code{exact} & 3 & 189/648 & 699 & 32.00 [18.12--34.90] \\
\pkg{RANN} exact & \code{exact} & 3 & 114/432 & 459 & 37.69 [22.16--39.21] \\
\pkg{Rnanoflann} exact & \code{exact} & 5 & 114/216 & 192 & 9.00 [8.60--21.35] \\
\pkg{BiocNeighbors} exact & \code{exact} & 7 & 238/432 & 286 & 3.16 [2.13--10.03] \\
\pkg{BiocNeighbors} HNSW & \code{hnsw} & 7 & 299/432 & 96 & 8.88 [4.08--10.71] \\
\pkg{RcppHNSW} HNSW & \code{hnsw} & 8 & 317/432 & 0 & 1.67 [0.77--2.95] \\
\addlinespace[0.4em]
\multicolumn{6}{@{}l}{\textit{Panel B: task-level alternatives}} \\
\addlinespace[0.25em]
\pkg{BiocNeighbors} Annoy & \code{auto} & 3 & 96/432 & 0 & 1.57 [1.46--1.73] \\
\pkg{RcppAnnoy} Annoy & \code{auto} & 2 & 89/432 & 6 & 2.15 [1.92--2.39] \\
\bottomrule
\end{tabular}
\caption{CPU comparisons on matched nodes. Panel A pairs exact and HNSW
interfaces with the corresponding \pkg{faissR} family. Panel B compares
Annoy interfaces with \pkg{faissR} automatic routing as a task-level
alternative, not as an Annoy implementation. Ratios are
$T_{\mathrm{comparator}}/T_{\mathrm{faissR}}$ among pairs in which both methods completed and met the requested mean recall. Each median and IQR is calculated from dataset medians; Datasets counts those with at least one eligible pair. Timeout counts sum route calls across both pair members and can exceed the pair count.}
\label{tab:comprehensive-r}
\end{table}

The results varied substantially across datasets and settings. These
conditional summaries describe the recorded public configurations and
supported metrics among successful pairs meeting the recall criterion. The
Annoy rows compare alternative interfaces for the same task and do not
measure an Annoy implementation in \pkg{faissR}. The separate HNSW experiment
provides the independently tuned same-family comparison.

\subsection{CUDA search and automatic selection}

The installed experimental policy selected exact-audited Flat in 280 of
324 combinations and IVF in 44. Supplementary
Tables~\ref{supp-tab:supp-auto-routes} and \ref{supp-tab:auto-ivf-parameters}
give the selections and IVF parameters. All combinations passed the recorded
validation criterion, but changing the requested tier changed the method
or parameters in only 10 of 108 dataset--metric--$k$ combinations.

We measured \emph{regret} as the selected method's runtime divided by that of
the fastest tested eligible method, including the selected method itself.
Median regret was 1.10, the 90th percentile was 3.38, and the maximum was
9.49. Supplementary Table~\ref{supp-tab:selector-regret-dataset} shows the
dataset summaries. IVF was faster than explicit exact search in 36 of its
44 selections, but the time lost on ImageNet exceeded the savings elsewhere.
The nine datasets are the primary experimental units; the 324 combinations
are repeated settings within them.

Aggregate and per-dataset CUDA ratios are reported in Supplementary
Tables~\ref{supp-tab:supp-cuda-paired} and \ref{supp-tab:supp-cuda-dataset}.
Grouped domain holdouts appear in
Supplementary Table~\ref{supp-tab:supp-domain-holdout}. These post hoc analyses
used exact methods and CAGRA, whereas the installed policy selected between
Flat and IVF. They assess sensitivity to the available methods and training
datasets; the preceding installed-policy analysis provides the corresponding
policy evaluation.

\section{Discussion and limitations} \label{sec:limitations}
\label{sec:conclusion}

\pkg{faissR} brings FAISS and optional cuVS search to R through a common
interface for dense input, reusable CPU indexes, and GPU-resident results.
The comparisons show why method choice should reflect both the data and the
query workload. In particular, the tuned HNSW results separate the cost of
constructing an index from that of querying it: a faster build does not imply
an equally large advantage for repeated queries.

The results describe the measured datasets and hardware. CPU pairs shared a
node, and the requested thread limits were recorded, while actual internal
thread utilization remained unmeasured. The broad seven-package experiment
used fixed settings. Its principal results separate same-family comparisons
from Annoy-versus-automatic-routing alternatives, while the experimental
derived-route comparison is confined to the supplement. The narrower HNSW
study tuned each implementation independently. Timeouts and failures are
reported alongside the ratios.

The optional automatic policy remains experimental. Its CUDA calibration
used an L40S and full self-search including construction, while CPU automatic
selection was not independently validated. The policy has no memory budget
or expected-reuse input. Shape descriptors capture size and dimension but
omit other aspects of data geometry, and one timing per calibration
configuration leaves selection sensitive to timing noise. The holdout analyses
cover the measured domains and hardware. Exact search remains a useful
baseline for application-specific comparisons beyond that scope.

GPU-resident output allows another compiled package to use neighbor buffers
without a mandatory host copy. The reported tests cover ownership and transfer
behavior; transfer speed remains unmeasured. Native indexes and GPU handles
remain local to the R session. Together, these facilities let users choose
among compiled search implementations while inspecting the method, input
conversion, and index reuse behind each result.

\section*{Data and code availability}

Package source, documentation, tests, manuscript source, and replication code
are available at \url{https://github.com/tkcaccia/faissR}. The replication
bundle provides checksummed benchmark outputs, analysis scripts, and a
machine-readable evidence manifest. Archive reanalysis requires no GPU; the
compact CPU mode executes the Biobase interface example. Raw execution
identities are retained in the manifests. Dataset sources, preprocessing,
fingerprints, and acquisition or access requirements are documented in the
supplement. Restricted data are not redistributed. The repository and
checksum ledgers identify the supplied material; a persistent archival
identifier is not currently available.

\section*{Acknowledgments}

The authors thank the University of Cape Town's Information and Communication
Technology Services (ICTS) High Performance
Computing team for providing a high-performance computing facility for this
study (\url{https://ucthpc.uct.ac.za/}). We also acknowledge the developers and
communities of FAISS, RAPIDS cuVS, CUDA, \proglang{R}, \pkg{Rcpp},
\pkg{float}, Singularity, and Bioconductor.

\bibliography{faissR_jss}

\end{document}
